\documentclass[10pt,a4paper]{article}
\usepackage[margin=20mm,headheight=15pt,headsep=7mm,footskip=10mm]{geometry}
\usepackage{fontspec,amsmath,amssymb,mathtools,microtype}
\setmainfont{lmroman10-regular.otf}[BoldFont=lmroman10-bold.otf,ItalicFont=lmroman10-italic.otf,BoldItalicFont=lmroman10-bolditalic.otf]
\setsansfont{lmsans10-regular.otf}[BoldFont=lmsans10-bold.otf,ItalicFont=lmsans10-oblique.otf,BoldItalicFont=lmsans10-boldoblique.otf]
\setmonofont[Scale=.86]{lmmono10-regular.otf}[ItalicFont=lmmono10-italic.otf]
\usepackage{ragged2e}\usepackage{xcolor,booktabs,longtable,array,graphicx,enumitem,fancyhdr,needspace,titlesec,hyperref,xurl}
\definecolor{AtlasNavy}{HTML}{18334C}\definecolor{AtlasTeal}{HTML}{116D73}\definecolor{AtlasGray}{HTML}{52616D}
\hypersetup{colorlinks=true,linkcolor=AtlasTeal,urlcolor=AtlasTeal,pdfcreator={Complexity Atlas},unicode=true}
\urlstyle{same}\setlength{\parindent}{0pt}\setlength{\parskip}{4.8pt}\setlength{\emergencystretch}{3em}
\renewcommand{\baselinestretch}{1.08}\setcounter{secnumdepth}{0}\setlist{nosep,leftmargin=1.5em}
\titleformat{\section}{\Large\bfseries\sffamily\color{AtlasNavy}\raggedright}{}{0em}{}
\titleformat{\subsection}{\large\bfseries\sffamily\color{AtlasTeal}\raggedright}{}{0em}{}
\titleformat{\subsubsection}{\normalsize\bfseries\color{AtlasNavy}\raggedright}{}{0em}{}
\titlespacing*{\section}{0pt}{14pt}{6pt}\titlespacing*{\subsection}{0pt}{10pt}{5pt}
\pagestyle{fancy}\fancyhf{}\fancyhead[L]{\small\sffamily\color{AtlasGray}COMPLEXITY ATLAS}
\fancyhead[R]{\small\sffamily\color{AtlasGray}ENGLISH EDITION}\fancyfoot[C]{\small\thepage}\renewcommand{\headrulewidth}{.3pt}
\newcommand{\AtlasMath}[1]{\par\begingroup\centering\setbox0=\hbox{$\displaystyle #1$}\ifdim\wd0>1.4\linewidth\typeout{ATLAS WIDE FORMULA: \the\wd0}\fi\ifdim\wd0>\linewidth\resizebox{\linewidth}{!}{\box0}\else\box0\fi\par\endgroup}
\newcommand{\AtlasInlineMath}[1]{\begingroup\setbox0=\hbox{$#1$}\ifdim\wd0>.4\linewidth\endgroup\AtlasMath{#1}\else\endgroup$#1$\fi}
\pdfstringdefDisableCommands{\def\AtlasInlineMath#1{#1}}
\DeclareMathOperator{\polylog}{polylog}\DeclareMathOperator{\dist}{dist}\DeclareMathOperator{\diam}{diam}
\DeclareMathOperator{\rank}{rank}\DeclareMathOperator{\tr}{tr}\DeclareMathOperator{\diag}{diag}
\DeclareMathOperator{\prox}{prox}\DeclareMathOperator{\KL}{KL}\DeclareMathOperator{\Var}{Var}
\providecommand{\eps}{\varepsilon}\providecommand{\R}{\mathbb R}\providecommand{\E}{\mathbb E}
\providecommand{\mathbbm}[1]{\mathbf{#1}}\providecommand{\mathds}[1]{\mathbf{#1}}
\providecommand{\llbracket}{[\![}\providecommand{\rrbracket}{]\!]}
\begin{document}
\raggedright

{\LARGE\bfseries\sffamily\color{AtlasNavy} Nonconvex-PL optimization: fresh gradients and weighted stationarity.\par}\vspace{5pt}
\textbf{Setting ID.} \texttt{CV\_\allowbreak{}S\_\allowbreak{}NCPL\_\allowbreak{}FRESH\_\allowbreak{}WEIGHTED}

Minimax optimization

Solve \AtlasInlineMath{\min_x\Phi(x)}, where \AtlasInlineMath{\Phi(x)=\max_yf(x,y)}.

\textbf{Status.} Bounds missing

\textbf{TASK tags.} solution concept: Weighted stationarity; task family: Minimax

\textbf{MODEL tags.} growth: Dual pl; initialization: Initialization; smoothness: Lipschitz

\textbf{ORACLE tags.} feedback quality: Stochastic unbiased bounded variance; interface family: Fresh joint gradient; operation: Stochastic joint gradient; public information: Hidden objects explicitly declared; public information: Public inputs explicitly declared; scalar derivative order: Scalar-function first derivative

\textbf{ALGORITHMS tags.} adaptivity: Adaptive; randomization: Private randomness allowed; restriction: Explicitly no additional algorithm restrictions; restriction: No additional span restriction explicitly imposed; restriction: No additional zero-respecting restriction explicitly imposed

\textbf{OUTPUT tags.} certificate: No additional output certificate required; object: Point

\textbf{METRIC tags.} criterion: Weighted joint gradient

\textbf{GUARANTEE tags.} aggregation: Second moment of the complete error; mode: Expectation; quantifier: Algorithm before instance; randomness: Algorithm private seed; randomness: Oracle samples

\textbf{RESOURCE tags.} accounting: Joint coordinates retained; budget: Fixed worst-case budget; measure: N grad; setup: Initialization included

\needspace{7\baselineskip}\section{1. Setting specification}
\needspace{4\baselineskip}\subsection{Task}
Solve \AtlasInlineMath{\min_x\Phi(x)}, where \AtlasInlineMath{\Phi(x)=\max_yf(x,y)}.

\needspace{4\baselineskip}\subsection{Admissible model}
The function \AtlasInlineMath{f:\mathbb{R}^{d_x+d_y}\to\mathbb{R}} is continuously differentiable with \AtlasInlineMath{\beta}-Lipschitz joint gradient. For each \AtlasInlineMath{x}, the value \AtlasInlineMath{\Phi(x)=\max_y f(x,y)} is finite and attained, and the dual PL condition holds:


\AtlasMath{\|\nabla_y f(x,y)\|_2^2\ge2\mu_y\bigl(\Phi(x)-f(x,y)\bigr).}


Concavity in the \AtlasInlineMath{y} direction and uniqueness of the maximizer are not required; set \AtlasInlineMath{\kappa=\beta/\mu_y\ge1}. Global PL and smoothness imply a differentiable value function, and \AtlasInlineMath{\nabla\Phi(x)} is independent of the selected maximizer.

Assume \AtlasInlineMath{\inf\Phi> -\infty} and \AtlasInlineMath{\Phi(x_0)-\inf\Phi\le H}. Define \AtlasInlineMath{h_{x_0}(x,y)=f(x,y)+\beta\|x-x_0\|_2^2}, and additionally require


\AtlasMath{b_0=2\bigl[\max_yh_{x_0}(x_0,y)-\min_xh_{x_0}(x,y_0)\bigr]\le c_{\rm init}H.}


The quantity \AtlasInlineMath{c_{\rm init}} is a fixed public constant. This initialization condition does not follow automatically from the primal gap.

\needspace{4\baselineskip}\subsection{Oracle interface}
Public inputs include block dimensions \AtlasInlineMath{d_x,d_y} and \AtlasInlineMath{H}, \AtlasInlineMath{\beta}, \AtlasInlineMath{c_{\rm init}}, \AtlasInlineMath{\varepsilon}, \AtlasInlineMath{\mu_y}, \AtlasInlineMath{\sigma}, together with the stated initial point and feasible geometry. The analytic representation of the hidden function, operator, or matrix is not revealed.


\AtlasMath{(x,y)\longmapsto G_f(x,y;\xi)=(G_x,G_y).}


Let \AtlasInlineMath{\mathcal H_t} denote the full history before query t. Conditional on \AtlasInlineMath{\mathcal H_t,z_t}, the response \AtlasInlineMath{G_t\in\mathbb R^{d_x+d_y}} satisfies \AtlasInlineMath{\mathbb E[G_t\mid\mathcal H_t,z_t]=\nabla f(z_t)} and \AtlasInlineMath{\mathbb E[\|G_t-\nabla f(z_t)\|_2^2\mid\mathcal H_t,z_t]\le\sigma^2}. Each call uses an independent fresh sample.

No reusable sample identifier or cross-point sample reuse is available. No mean-square smoothness of sample gradients or operators is additionally assumed. Blocks within a full vector may be correlated.

No best response, reduced gradient, or function value is supplied.

\needspace{4\baselineskip}\subsection{Output}
Return \AtlasInlineMath{(\widehat{x},\widehat{y})\in\mathbb{R}^{d_x+d_y}}.

The output may be any measurable postprocessing of the history, including an unqueried point, average, or convex combination. It must satisfy the stated domain and feasibility requirements; no additional function value or computable certificate is required.

\needspace{4\baselineskip}\subsection{Loss or accuracy criterion}

\AtlasMath{\bigl(\|\nabla_x f(\widehat{x},\widehat{y})\|_2^2+\kappa\|\nabla_y f(\widehat{x},\widehat{y})\|_2^2\bigr)^{1/2}}


\needspace{4\baselineskip}\subsection{Resource accounting}
One complete joint-gradient sample at one point costs one call. Using independent samples at two points in an iteration costs two.

Use a fixed worst-case budget: calls never exceed the prespecified \AtlasInlineMath{N} on any instance or random execution path. Earlier stopping within the budget is permitted.

Count all information calls for one instance and one requested accuracy, including initialization. Exact-real arithmetic, storage, and solving known models from public inputs and returned data are uncharged. This query-complexity quantity does not certify arithmetic or bit complexity.

\needspace{4\baselineskip}\subsection{Algorithm views and success guarantees}
Arbitrary measurable sequential adaptive randomized algorithms are allowed. The private seed is independent of the hidden instance. Queries, stopping, and output may use the full history, without additional span, zero-respecting, communication, or memory restrictions.

Choose one algorithm from public inputs, then require the guarantee for every admissible fixed instance and fresh-sample feedback distribution.


\AtlasMath{\mathbb E\!\left[\|\nabla_x f(\widehat x,\widehat y)\|_2^2+\kappa\|\nabla_y f(\widehat x,\widehat y)\|_2^2\right]\le\varepsilon^2.}


Expectation covers uniform output selection and all independent fresh samples.

\textbf{Budget convention.} fixed\_worst\_case\_cap

\needspace{7\baselineskip}\section{2. Notation and parameters}
\AtlasInlineMath{H}: Public scale jointly controlling the initial reduced-objective gap and auxiliary initial saddle gap \textbf{Domain:} Positive real number.

\AtlasInlineMath{\Phi(x_0)-\inf\Phi\le H,\qquad b_0\le c_{\rm init}H}

\AtlasInlineMath{\beta}: Euclidean Lipschitz upper bound for the population payoff's complete joint gradient. \textbf{Domain:} Positive real number.

\AtlasInlineMath{\|\nabla f(z)-\nabla f(w)\|_2\le\beta\|z-w\|_2}

\AtlasInlineMath{c_{\rm init}}: Fixed upper bound on the ratio between an auxiliary initial gap and \AtlasInlineMath{H}. \textbf{Domain:} A fixed positive constant.

\AtlasInlineMath{d_x}: Primal-space dimension \textbf{Domain:} Positive integer.

\AtlasInlineMath{d_y}: The dimension of the dual-variable space. \textbf{Domain:} Positive integer.

\AtlasInlineMath{\varepsilon}: Root-mean-square tolerance for the joint residual comprising the primal gradient and the \AtlasInlineMath{\kappa}-weighted dual gradient. \textbf{Domain:} Positive real number.

\AtlasInlineMath{\kappa}: Ratio of the joint smoothness bound to the dual PL constant \textbf{Domain:} \AtlasInlineMath{\kappa\ge1}.

\AtlasInlineMath{\kappa=\beta/\mu_y}

\AtlasInlineMath{\mu_y}: Uniform positive constant in the dual PL inequality for every fixed x \textbf{Domain:} \AtlasInlineMath{0<\mu_y\le\beta}.

\AtlasInlineMath{\|\nabla_yf(x,y)\|_2^2\ge2\mu_y(\Phi(x)-f(x,y))}

\AtlasInlineMath{\sigma}: Standard-deviation upper bound for noise in the full stochastic response, controlling the joint squared norm rather than summing independent block budgets. \textbf{Domain:} A nonnegative real.

\AtlasInlineMath{\mathbb E[\|G_t-\mathbb E[G_t\mid\mathcal H_t,z_t]\|_2^2\mid\mathcal H_t,z_t]\le\sigma^2}

\AtlasInlineMath{\|\cdot\|_2,\langle\cdot,\cdot\rangle}: Euclidean norm and inner product; the product-space norm is the square root of the sum of the squared block norms. \textbf{Domain:} The specified finite-dimensional real vector space.

\AtlasInlineMath{\|v\|_2=\sqrt{\langle v,v\rangle}}

\AtlasInlineMath{f}: Unknown real-valued payoff, with primal and dual input blocks. \textbf{Domain:} \AtlasInlineMath{f:(x,y)\mapsto f(x,y)\in\mathbb R}.

\AtlasInlineMath{x,y,z}: Primal variable, dual variable, and joint point \textbf{Domain:} \AtlasInlineMath{x\in\mathbb R^{d_x},\ y\in\mathbb R^{d_y},\ z=(x,y)\in\mathbb R^{d_x+d_y}}.

\AtlasInlineMath{z=(x,y)}

\AtlasInlineMath{\nabla f,\nabla_x f,\nabla_y f}: Complete joint gradient and its two blocks \textbf{Domain:} \AtlasInlineMath{\nabla_x f\in\mathbb R^{d_x},\ \nabla_y f\in\mathbb R^{d_y}}.

\AtlasInlineMath{\nabla f=(\nabla_x f,\nabla_y f)}

\AtlasInlineMath{\Phi}: Reduced objective after global dual maximization; neither its value nor gradient is an additional free interface \textbf{Domain:} Real-valued function on the primal-variable space.

\AtlasInlineMath{\Phi(x)=\max_{y\in \mathbb R^{d_y}}f(x,y)}

\AtlasInlineMath{\operatorname{Argmax}_y f(x,y)}: The nonempty set of global dual maximizers at fixed x. It need not be a singleton. \textbf{Domain:} Nonempty closed subset of the dual-variable space.

\AtlasInlineMath{\{y:f(x,y)=\Phi(x)\}}

\AtlasInlineMath{x_0,y_0}: The public initial primal and dual variables. \textbf{Domain:} The respective feasible domains. Both equal zero when the origins are explicitly fixed.

\AtlasInlineMath{\widehat x,\widehat y}: The point or variable blocks ultimately returned by the algorithm, corresponding to the hatted variables in the error formula. \textbf{Domain:} Feasible variable pairs in \AtlasInlineMath{\mathbb R^{d_x}\times\mathbb R^{d_y}}.

\AtlasInlineMath{G_f(z;\xi),G_t}: Random sample of the full joint gradient and the response from sample \AtlasInlineMath{t} \textbf{Domain:} \AtlasInlineMath{\mathbb R^{d_x+d_y}}. The two blocks may be internally correlated.

\AtlasInlineMath{G_f=(G_x,G_y)}

\AtlasInlineMath{G_x,G_y}: The primal and dual blocks in one complete joint-gradient sample. \textbf{Domain:} \AtlasInlineMath{G_x\in\mathbb R^{d_x},\ G_y\in\mathbb R^{d_y}}.

\AtlasInlineMath{\xi,\mathcal H_t,z_t}: Current fresh-sample seed, complete pre-query history, and query point. \textbf{Domain:} The seed space, history-generated \AtlasInlineMath{\sigma}-algebra, and declared query domain.

\AtlasInlineMath{\mathbb E,\Pr}: Expectation and probability over the explicitly allowed algorithm and sample randomness, with the hidden instance fixed. \textbf{Domain:} Operations on the corresponding probability space.

\AtlasInlineMath{N_{\rm grad}}: Number of full joint-gradient queries using fresh samples \textbf{Domain:} Nonnegative integer count; expected work is its expectation.

\AtlasInlineMath{N}: A uniform cap on principal information calls for every admissible instance and execution path. \textbf{Domain:} Nonnegative integer.

\AtlasInlineMath{N^\star}: The infimum principal query budget among algorithms satisfying every stated instancewise output guarantee. \textbf{Domain:} Nonnegative extended real value; by convention, positive infinity if no successful algorithm or no uniform finite budget exists.

\AtlasInlineMath{\mathcal A,I}: An algorithm selected from public inputs, a hidden admissible object fixed before interaction, and an oracle law satisfying the declared conditional means, variances, and independence protocol. \textbf{Domain:} The specified algorithm and model classes.

\AtlasInlineMath{\omega}: A complete realization of algorithmic and sample randomness. \textbf{Domain:} Corresponding sample space.

\AtlasInlineMath{O,\Omega,\Theta,\widetilde O}: Asymptotic upper and lower bounds, constant-factor matching, and upper bounds hiding the stated logarithmic factors, within the corresponding regions. \textbf{Domain:} Nonnegative complexity function; subscript p permits additional dependence only on fixed p.

\AtlasInlineMath{h_{x_0}}: Regularized payoff centered at the public initial primal variable \textbf{Domain:} A real-valued function on \AtlasInlineMath{\mathbb R^{d_x+d_y}}.

\AtlasInlineMath{h_{x_0}(x,y)=f(x,y)+\beta\|x-x_0\|_2^2}

\AtlasInlineMath{b_0}: Twice the primal-dual gap of the first regularized saddle problem at the public initial point \textbf{Domain:} A nonnegative real.

\AtlasInlineMath{b_0=2[\max_y h_{x_0}(x_0,y)-\min_x h_{x_0}(x,y_0)]}

\needspace{12\baselineskip}\section{3. Complexity bracket}
\needspace{4\baselineskip}\subsection{Compared quantities}
\needspace{4\baselineskip}\subsubsection{Q1 - source scoped metric resource statement}
\textbf{Metric contract.} output: domain: R\textasciicircum{}\{d\_x\} x R\textasciicircum{}\{d\_y\}; postprocessing: arbitrary measurable transcript-based output unless algorithm restriction expressly states otherwise; metric: sqrt(||grad\_x f(xhat,yhat)||\textasciicircum{}2+kappa||grad\_y f(xhat,yhat)||\textasciicircum{}2); guarantee: One algorithm family chosen from public inputs satisfies for every legal instance and oracle law: E[(sqrt(||grad\_x f(xhat,yhat)||\textasciicircum{}2+kappa||grad\_y f(xhat,yhat)||\textasciicircum{}2))\textasciicircum{}2]<=epsilon\textasciicircum{}2.

\textbf{Resource contract.} budget: fixed worst-case cap on every transcript; primary: N\_grad; unit: one declared oracle call; auxiliary: projection/internal model solves separately retained, never counted as hidden oracle

This is the native source metric/resource guarantee. Joint precision, block size, runtime and oracle components stay separate; do not infer a scalar optimal query complexity.

\needspace{4\baselineskip}\subsection{R1 - The bound’s native applicability window.}
\textbf{Region.} Assume \AtlasInlineMath{\varepsilon>0}, \AtlasInlineMath{\Phi(x_0)-\inf\Phi\le H}, and \AtlasInlineMath{0\le b_0\le c_{\rm init}H}, with fixed public constant \AtlasInlineMath{c_{\rm init}}. All remaining smoothness, dual PL, and fresh-sample noise conditions are those specified in this profile.

\textbf{Quantity.} Q1

\textbf{Resource.} Joint resource budget

\textbf{Status.} Lower bound missing

\textbf{Lower bound.} No admitted bound for this comparison

\textbf{Upper bound.} B1

\needspace{7\baselineskip}\section{4. Bounds and supporting evidence}
The following conditions apply to every bound in this profile, in addition to each bound's stated region and algorithm scope.

\textbf{Conditions.} Equations (40)-(41) control the weighted squared residuals of both blocks, with initialization \AtlasInlineMath{\Delta+b_0\le(1+c_{\rm init})H}. Suitable step sizes and random iterate selection give the stated mean-square guarantee. Simulate the two independent block samples with two joint-sample calls, discard the unused blocks, and charge both calls.

\needspace{4\baselineskip}\subsection{B1 - See the cited theorem}
\textbf{Admission.} Included in the bracket

\textbf{Direction.} upper

\textbf{Witness or lower-bound scope.} Randomized witness

\textbf{Algorithm-scope qualification.} Stochastic Smoothed-AGDA in source Algorithm 2 uses random samples and draws its output uniformly from the iterate history.

\AtlasMath{O\!\left(1+\frac{\beta\kappa H}{\varepsilon^2}+\frac{\beta\kappa^2 H\sigma^2}{\varepsilon^4}\right)}
\textbf{Sources.} \href{https://proceedings.mlr.press/v151/yang22b/yang22b.pdf}{[1]}

\textbf{Constants.} Use the source theorem’s universal constants except where derivative order or structural budgets are explicitly fixed. Unspecified sufficient thresholds remain symbolic.

\textbf{Source location.} \href{https://proceedings.mlr.press/v151/yang22b/yang22b.pdf}{[1]} Assumptions2.1-2.3; Algorithm2; Theorem4.1; Remark4.1; proof Eq40-41

Fresh-sample Smoothed-AGDA controls weighted squared joint gradient.

Block Lipschitz source constant l; beta-Euclidean joint smoothness implies source assumptions with l=beta. Source per-block variance <=sigma\textasciicircum{}2 follows from joint variance <=sigma\textasciicircum{}2. Two independent sample calls each iteration at different points. Explicit auxiliary initialization b0 must be retained, or b0<=cH and primal gap<=H imposed. Weighted RMS target is not direct primal gradient.

\needspace{7\baselineskip}\section{References}
[1] Junchi Yang; Antonio Orvieto; Aurelien Lucchi; Niao He. Faster Single-loop Algorithms for Minimax Optimization without Strong Concavity. 2022. AISTATS2022 PMLR151.

\url{https://proceedings.mlr.press/v151/yang22b/yang22b.pdf}

\end{document}
